\chapter{实现}
\label{chap:implementation}

我们基于 DPDK 实现 \sysname 主机运行时，并为 Alveo V80 实现四端口聚合
RTL kernel。物理原型通过 Tofino 连接四个 100-GbE endpoint 端口和四个
V80 accelerator 端口；交换机只执行 unicast/multicast 转发。V80 kernel
暴露四组 packet AXI4-Stream 接口，DCMAC/GTM shell 负责物理网口、FEC、
GT reset 和时钟转换。六种集合通信原语共用同一 channel 数据路径，原语只
决定 channel 集合和 payload 方向。

\section{原型结构}
\label{sec:impl-structure}

\parab{Channel 单写与 shared poller.}
主机运行时由 EAL lcore poller 构成。每条 physical channel 拥有一对 DPDK
队列及独立的 credit、packet-offset 和 message 状态；配置把每条活跃 channel
恰好分配给一个 poller，一个 poller 可以轮询多条 channel。subchannel 使用
不同 UDP 端口，NIC flow steering 将其导入所属 poller 的 RX 队列，poller
再按报头分派到 channel。该映射保持 channel 状态单写，多条 channel 合并轮询
时不引入跨核锁。当前原型为每个本地运行时实例分配独占的 DPDK device 或 VF；
不同通信组仍共享网络中的 allocator。

\parab{异步 message 发布.}
一次 primitive 调用产生若干 channel message，每个 message 再切成
8192-byte packet。异步接口先在固定表中建立全部 message 并预留 packet
offset，随后用 256-bit 原子位图一次性发布；准备失败的 collective 不进入
数据面。调用在发布后立即返回，worker 根据位图事件启动 per-message 注册。
该路径只串行化同一 communicator 的提交，已发布 message 可以并行推进。

\parab{报头与路径编码.}
多级 profile 使用 20-byte header；物理 testbed 的单级 V3 profile 是前
16 bytes 的严格截断，省略后两级位置字段。两种 profile 的前 16 bytes 具有
相同字段和偏移，分别携带 payload offset 和 credit offset，并用独立
valid bit 区分返回数据与发送许可。24-bit packet offset 在 channel 内连续
编号；UDP port 编码 subchannel，因此 slot owner 可以表示为
\texttt{(UDP port, packet offset)}。Normal credit 最多携带三级位置栈，
2-bit \texttt{agg\_depth} 给出栈深，request 用三个对齐字节携带逐级
fan-in。控制包用 message 起始 offset 作为标识符复用 epoch。Normal
credit 只有在取得沿途全部位置后才交付 requester；\texttt{AGG\_MISS}
清除 payload 和位置栈。

\parab{NUMA 局部性与并行扩展.}
$R_{\rm rep}$ 个 replica 为一个 logical channel 建立相互独立的 physical
channel；poller 数量则决定这些 channel 使用多少主机核。两者相互独立：多条
channel 可以共享一个 poller，增加 poller 也不改变 channel、subchannel 或
credit 的协议语义。每个 channel 的 NIC 队列、mbuf pool 和状态环按其 poller
放置到本地 NUMA 节点；跨 socket 绑定需要显式启用。用户 buffer 的放置仍由
上层应用决定。\Cref{sssec:eval-endpoint-cpu} 分别扫描 replica 与 poller，
并报告达到目标链路速率所需的最小协议核数。

\parab{零拷贝 request 与分段接收.}
运行时把用户输入 buffer 注册给 NIC，并用 DPDK external buffer 将
payload 切片直接挂到 header mbuf 后。首发和 repair request 共用该
payload；repair 只重建 header。接收路径使用 NIC RX buffer split 分离
header 与 payload。requester 解析 header 后，将 response payload 一次
拷贝到用户输出 buffer。

\parab{可选的 I/OAT 交付路径.}
写回用户输出 buffer 的这次搬运有两条已注册路径：poller 直接执行 CPU
拷贝，或交给同 NUMA 节点的 I/OAT DMA 引擎、poller 只提交描述符并回收
完成。两条路径不改变协议语义，也不增加 poller；按 primitive 和消息
大小的路径选择在正式测量前冻结
（\cref{app:physical-implementation-audit}）。

\parab{Requester 无定时器、无重传状态.}
requester 维护一个 credit 队列和一个按 packet offset 取模的去重环，环深
为 $2\times\texttt{MAX\_WINDOW\_SIZE}$。Normal offset 单调推进，稳态
message 查找只检查当前与下一条 message；repair 才扫描完整表。Response
payload 按 \texttt{payload\_offset} 去重写回，credit 则连同 subchannel、
位置栈和 CE 状态入队。发送循环逐项构造 request，并带回位置栈、fan-in
和 CE。

每个未确认 message 的 \texttt{REGISTER} 在 channel 级预算内轮转发送；
\texttt{REGISTER\_ACK} 或首份 credit 停止其信标。注册门阻止 responder 在
全组 message 状态发布前发放 normal credit。requester 不维护数据重传
定时器；收到 repair trigger 后重发 repair request，由 responder 位图
去重（\cref{ssec:reliability}）。

\parab{未完成 credit 表与聚合 slot 合一.}
responder 将未完成 credit 表与主机聚合状态合并为一个按 packet offset
取模的状态环，深度同样为 $2\times\texttt{MAX\_WINDOW\_SIZE}$。entry 在
credit 发放时创建，并保存 requester 位图、master response、subchannel、
提交/repair 状态、重试次数、时间戳和拥塞控制参考状态。端侧拥塞控制模块
取得首个有效样本前，启动窗 $W_0$ 限制总在途 credit。环绕遇到仍承担重放
责任的 entry 时，responder 暂停发放，不会覆盖未确认结果。

交换机聚合结果和 repair 路径的主机聚合结果写入同一份 master response。
已提交结果用 mbuf clone 重放，不重新执行聚合。worker 每轮按固定步长扫描
状态环；超时触发组播 repair，独立 token bucket 限制 repair 流量。Master
response 保留到绑定它的后续 credit 完成；没有后续 credit 的尾部结果保留
到 \texttt{END\_ACK}。

\parab{消息结束与标识符复用.}
\texttt{END}、\texttt{END\_ACK} 和注册包都携带 message 起始 offset 作为
复用 epoch。requester 为最近完成的 epoch 保留 tombstone，因此 message
ID 复用后仍能回答迟到的 \texttt{END}。responder 在收齐 \texttt{END\_ACK}
前保留尾部 response；状态环绕回这些 entry 时暂停新 credit。

\section{聚合加速器：协议把数据面化简为计数与累加}
\label{sec:impl-accel}

端点维护单注入、成员和恢复状态；加速器只保存聚合中的 Credit Context。
加速器旁接商用交换机，不安装通信组、rank、IP/MAC 或转发表。

\parab{Request 自含聚合参数.}
request 用 credit 返回的位置栈直接索引本级 slot，并用
\texttt{(UDP port, packet offset)} 验证 owner。fan-in、算子和数据类型均
随 request 到达。RTL 解析固定偏移字段、推进位置栈，并完成 owner 比较、
计数和累加。当前 data-type profile 支持 FP32 的 SUM/MAX/MIN；其它数值
格式是 profile 扩展，不改变协议字段或恢复规则。聚合结果携带贡献计数，
responder 提交前验证其等于 requester group 规模；可选 ECN 路径在同一
slot 内对 CE 取 OR。控制包和 repair request 不进入加速器，拥塞控制状态
位于 responder。

\parab{预装 channel 状态跨 message 复用.}
当前原型的主机运行时持续处理所有预装 channel。Endpoint 创建 communicator
时只建立本地 channel 状态并完成 \texttt{REGISTER}；销毁前等待
\texttt{END}/\texttt{END\_ACK} 收尾，随后释放本地状态。该过程不重启 EAL，
也不向 V80 写入 per-group 配置。网络侧只需在普通交换机中预装常规单播和
组播转发表项；Arbor 不依赖可编程交换机。当前原型在建组前固定成员关系，
运行中不更新这些表项。

测试床中每个 subchannel 经过一个聚合阶段，位置栈退化为单项。包级仿真
实现并评估多级语义（\cref{ssec:agg-feedback}）。

\parab{分 bank 的 slot 池.}
所有物理端口共享一个 allocator。RTL 参考配置使用 1024-slot 池，并按
slot 低位把状态交织到四个 bank；每个 bank 包含 256 项元数据、payload
memory 和 FP32 pipeline。通用 RTL core 的 packet port 数由
\texttt{NUM\_PORTS} 参数化；V80 kernel 固定为四个 packet ports，bank 数
固定为四。入口使用 $\texttt{NUM\_PORTS}\times4$ 个 VOQ 按目标 bank 排队，
出口使用 $4\times\texttt{NUM\_PORTS}$ 个 VOQ 返回原 port。四个 worker
可以同时处理不同 Credit Context；同一 slot 的 fan-in 由一个 bank 顺序
累加，无 payload request 只更新贡献计数。

全局 allocator 为每个 normal credit 生成唯一 ticket；ticket $n$ 直接映射
到 $n\bmod A_e$，其低位同时选择 bank。四个 bank 只是同一分配序列的并行
实现，不建立独立指针或 slot 池。池深按 \cref{sssec:slot-bound} 的循环
复用条件配置；迟到 request 的 owner mismatch 生成 \texttt{AGG\_MISS}，
冲突计数用于检查配置。

\parab{聚合输出与组播注入分离.}
本级 request 聚合完成后，加速器只沿 Credit Context 绑定的上行父端口返回
一份结果；responder 收到结果后，才从该 subchannel 发出下一份
response/credit 并由普通交换机组播。四个 bank 可以并行完成不同 Context，
但这些完成流不会分别复制到同一组 requester。每个 subchannel 的速率门和
窗口门同时限制 responder 的组播注入，因此加速器不需要在四个物理端口之间
维护一份面向所有 requester 的全局结果发送队列。

该路径与结果本身需要组播的固定聚合接口不同。SwitchML~\cite{switchml}
若把一个逻辑聚合器分布到四条 FPGA 链路，四条独立完成流在进入交换机组播
复制后都可能同时到达每个 worker egress；长期的四输入归约为一输出并不能
限制这一瞬时峰值。这样的实现需要在组播前串行化全局完成流，并缓存等待发送
的完成描述符。ATP~\cite{atp} 的 worker aggregate 则先单播汇聚到 PS，PS
产生的 parameter 构成单一组播输入，因而没有相同的跨端口复制峰值；它仍需
维护 hash owner、collision/remap 和 parameter GC。Arbor 将位置、owner、
fan-in 和恢复边界放入 credit/request 闭环，使加速器的协议控制状态限于
直接索引的 owner、贡献计数和完成描述符。资源比较据此将 payload accumulator
与这些协议控制资源分开报告。

\parab{直接索引的 slot 数据路径.}
数据面状态是一个直接索引 slot 数组。Normal request 单注入使 slot 可以按
贡献数推进；repair request 不进入加速器（\cref{ssec:reliability}）。端点
按模 $2^{24}$ 比较 packet offset，并让存活 message 保持在无歧义半空间
内。\texttt{END} 保护 message ID 复用，未完成 credit 表和 owner 校验共同
保护在途 DATA。由此，RTL 无需 CAM、哈希表、per-flow 速率状态、
per-requester 记录或逐包定时器。

\Cref{chap:evaluation} 按机制报告物理协议路径、FP32 正确性和包级仿真的
证据范围，\cref{chap:llm-evaluation} 报告大模型训练评估。
